\newenvironment{abstract}% 
{\thispagestyle{empty} \cleardoublepage\null
\thispagestyle{empty} 
\vfill\begin{center}% 
    \bfseries 
    \abstractname
\end{center}} 
{\thispagestyle{empty} \vfill\null }

% \selectlanguage{spanish} 
\begin{abstract}
Durante el desarrollo de esta Tesis se desarrolló en primera instancia un sensor capacitivo puntual para la medición de la altura de la superficie libre de agua a una alta tasa de adquisición (hasta algunos kHz) y alta precisión (resolución submilimétrica), que fue caracterizado y mostró una respuesta lineal en el rango de interés. Para esto fue diseñada la electrónica y el sistema de adquisición, compuesto por un lock-in digital que utiliza una placa de adquisición Personal DAQ 3001 de IOTech, para la cual también se desarrolló el \textit{software} pertinente. En una segunda instancia se utilizó dicho sensor en experiencias típicas de turbulencia de ondas, inyectando mediante un forzado aleatorio del fluido limitado en frecuencia, pudiendo replicar distintos resultados de la literatura en el caso estacionario, como por ejemplo los exponentes de los espectros de energía de Kolmogorov-Zakharov-Filonenko (KZF) en los rangos de gravedad y capilaridad,  y permitiendo además identificar y caracterizar distintos mecanismos para la transferencia de energía entre ondas, como lo son las tríadas y cuartetos resonantes. Además fue posible observar y estudiar para el caso de turbulencia con no linealidades altas el fenómeno de intermitencia. % sísísís aá a as 


\end{abstract}

\selectlanguage{english}  % 
\begin{abstract}  %


  		
During the development of this Thesis, a point-like capacitive sensor was first designed and implemented to measure the free surface elevation of water at a high acquisition rate (up to several kHz) and with high precision (submillimetric resolution). The sensor was characterized and exhibited a linear response within the range of interest. To this end, both the electronic circuitry and the acquisition system were developed, consisting of a digital lock-in based on a Personal DAQ 3001 acquisition board (IOTech), for which the corresponding software was also implemented.  	In a second stage, the sensor was employed in typical water wave turbulence experiments, where energy was injected into the fluid through random forcing within a limited frequency band. This allowed the reproductions of several results reported in the literature for the stationary regime, such as the Kolmogorov–Zakharov-Filonenko (KZF) energy spectrum exponents in both the gravity and capillary ranges. Furthermore, it enabled the identification and characterization of different mechanisms responsible for energy transfer between waves, such as resonant triads and quartets. Additionally, in the regime of strongly nonlinear turbulence, the phenomenon of intermittency was observed and analyzed. 
 	
 	 
 	
\end{abstract} % 
\selectlanguage{spanish} % 

